\chapter{研究背景}


\section{数学习题课的教学模式现状及存在问题}


\subsection{教学模式不科学，不注重问题引导，轻过程重结论}

以目前的高中课堂看来，习题课多以教师为主导地位，具体表现为教师课上带领学生进行数学习题训练，或是学生课下提前做题，课上教师讲评的方式。并且，在习题讲解过程中，教师先入为主，一股脑地将答案呈现在黑板，没有过程的思考与推导。这种模式古板单一，学生更多的是被动地接受知识和解题经验的传输，教师习惯于一昧传输自身所学，忽视了学生该有的思考，抛弃了课标中所强调的“自主探究”的学习方式。


\subsection{选题量大盲目且缺少衍生性，摒弃培养学生的创新思维}

虽然新课改已实施近20年，但“应试教育”的观念仍然积重难返，其中特别明显地体现在习题课的教学上。教师依然以高考为终极目标，围绕高考的命制原则进行习题课内容的选择，抛弃了与考试无关的习题，例如以数学文化为背景的长题型、开放性问题及探究性问题等。但是这些问题恰恰是课标中所强调的，命题原则里明确指出，题目应包含具有探究性的数学问题，并且需要有开放性问题，考查学生的思维能力。由于题海战术的深刻影响，一个很小的知识点需要三或四道题目进行巩固，因而教师选择的习题数量较大，看到好题、新题随手拿到课堂来分享，可见对于题目的选择没有计划与原则，更没有考虑到每一道习题的内核，即习题背后待发掘的原理。且教师常把每一道习题作为独立的个体，讲评完上一道跳到下一道，并没有用知识把习题之间串起来，对于题目的选择没有连贯思维，缺乏了衍生性，压抑了学生的创新思维与创新能力。


\subsection{习题课教学设计的重视性小}

许多教师对习题课的重视程度不够，对于习题课作用的看法只停留在解决平时作业、考试等练习中存在的错误或疑问。因此，教师对新授课倾向于花大量时间进行备课、教学设计，习题课与之相比，重视程度呈断崖式下降。更有不少老师认为，习题课是最轻松的课，只需要讲评题目就行了，根本不需要进行完整的教学设计。但在实际课堂上，由于先前教学设计不到位，短短45分钟的课堂上很容易出现习题多且不精、讲解不完全、时间把握失控等问题。



\section{数学习题课的学生学习现状及存在问题}

学生在习题课上的学习效率部分依赖于教师教学，因此由于教学模式存在问题，学生学习也存在相应的问题，例如兴致缺缺、缺乏思考、被动学习、效率低下等问题。除此之外，学生本身对习题课不够重视，从而导致听讲不认真、不主动，进而造成课堂效率低下的问题。学生会惯性认为教师在习题课上无非是公布答案、讲解习题解答过程，但由于现在网络的发达，学生可以在电子产品实时搜索到题目答案与解题过程，因此有部分学生会认为习题课没有听课的必要，并且对习题课提不起兴趣。



\section{习题课开展探究性学习的意义}

在《普通高中数学课程标准》的2020年修订版中提到，课程内容突出的主线之一是数学建模活动与数学探究活动\cite{普通高中数学课程标准}。数学探究是一种以数学知识为基础，以灵活、开放的方式来解决数学问题的一种具有思维高度综合活动，而探究性学习正是开展此活动的一个最为主要的学习方法。探究性学习是学生在理解本源知识的基础上，对系列问题进行开放性思考，提出问题并解决问题，最终获得知识和方法的一种主动性探究学习。因此，探究性学习在高中数学领域是被提倡的，被广泛应用的，且具有积极意义的学习方法之一。

习题课是高中数学里的重要课型之一，教师可以通过习题课的教学，检查学生对于已经教授的知识是否掌握，发现学生的知识漏洞，进而调整自己的教学方法和步骤，是教学活动中极为重要的实践反馈环节。同时，学生通过习题课也能够巩固所学的理论知识，促进解题能力的提升。课标强调高中课程目标包括“四能”，即从数学的角度发现和提出问题的能力、分析和解决问题的能力\cite{普通高中数学课程标准}。四能点明了高中与义务教育课程的区别，高中更加注重数学思维和核心素养的培养，习题课作为高中重要课型之一，也同样要强调思维和素养的提升。同时，“四能”也恰好对应了探究性学习的显著特征。在探究性学习中，学生一定会涉及到“做题”、“想题”、“解题”等和题目相关的过程，因此，探究性学习与题目密不可分，也就是说探究性学习和习题课之间也有着不少共通之处。在这种联系之下，我们可以把习题课与探究性学习相结合，探求习题课中的探究性学习到底是什么，教师如何教学，学生如何学习等相关问题，且这些问题是具有必要性和重要性的。

探究性学习提倡学生自主学习获得知识，深究知识，并用知识解决问题，更加关注学习过程，注重问题的发现，能够锻炼学生的开放性思维与创新能力。探究性学习等同于经历科学家发现科学定律的部分经历，在古往今来的数学发现中，几乎所有的发现都要经历“观察—猜想—证明”这些步骤，这些过程就是探究性学习的核心。因此，教师应当在数学教学中使用探究性学习，发挥学生的主观能动性，利用习题的本质特点延伸出有待探究的问题，引导学生开展探究性学习，使其体验到自主学习和合作交流的积极学习模式，感受到知识是靠自己习得的，拥有实现自我的满足感，从而体会到学习数学真正的乐趣。



\section{研究内容与主要研究方法}

在以往的研究中，大部分学者将探究性学习运用到高中数学的课堂教学中，运用在习题课中的案例少之又少。也有部分学者将习题与探究性学习联系起来，针对某个或一系列的习题提出具有探究性学习性质的解题方法，鲜少提到在习题课上教师应该如何开展探究性学习以及学生如何进行探究性学习。本文将探究性学习运用到习题课的教学中去，为传统习题课教学提供新的理论基础与实践方案。本文首先阐述了高中数学习题课的教学模式的存在问题以及进行习题课中的探究性学习的意义，评述国内外的研究现状，进行探究性学习等关键词的概念界定，结合选定的课题《函数的性质——函数的凹凸性》进行探究性学习设计方案，并在选定班级进行实验后进行反馈分析得到结论与不足。主要研究方法如下：


\subsection{文献研究法}

本文的研究立足于大量的文献成果，总结整理国内外学者们关于高中数学习题课中的探究性学习的研究，通过分析定义习题课中的探究性学习，提出关于《函数的性质——函数的凹凸性》的探究性学习具体实践方案，以普通高中课程标准为引导，在此基础上创新思考，提出新的研究方向。


\subsection{问卷调查法}

本文拟采用问卷调查的方法针对本文提出的关于《函数的性质——函数的凹凸性》的探究性学习具体实践方案，对本课堂的学生进行课堂评价、教学效果等多方面的问卷调查，通过线下进行问卷收集，最终整理后的有效问卷为数据分析提供支撑，根据数据直观给出该实例应用的评价与成效，进一步说明探究性学习在习题课中的应用。


\subsection{深层访谈法}

针对个别旁听教师进行深层开放式访谈，例如对于课堂的探究性学习有何想法、对该习题的认识如何等等，旨在分析实例应用的旁听教师的建议，进一步分析探究性学习在习题课中的优势与不足。


\subsection{综合分析法}

本文对所收集的数据进行整理、分析和数理统计，对问卷进行调查对象频率分析、信度检验，再针对数据直观给出统计描述性分析，调查实例应用的评价与成效，进一步说明探究性学习在习题课中发展的利弊，据此制定相应建议与改善。